\section{Related work}\label{sec:related}
% Outline row: related.tex. Autonomy frameworks; onboard decision support; wearables in flight; early warning
% from wearables; alarm design; operations under delay. Table 1 holds only systems whose sources were verified.

\paragraph{Autonomy frameworks.}
The idea that medical authority should move toward the crew as Earth gets further away predates this work. NASA's Exploration Medical Capability element defines EIMO as a process in which terrestrial assets keep pre-mission screening, planning and prevention while onboard care increasingly becomes the crew's responsibility~\cite{levin2023eimo}, and a NASA workshop report sets out the same transition for Lunar and then Martian missions~\cite{lemery2024eimo}. The Canadian Space Agency's expert group on deep-space healthcare framed Canada's possible role around autonomous care in space and on Earth~\cite{csaDeepSpace2019}. Cypko et al.\ propose four autonomy levels tied to mission distance and communication latency, from Earth orbit to deep space~\cite{cypko2025}. These works describe the levels. LIGHT-TIME computes which level applies from the current light time and shows it on every screen.

\paragraph{Onboard decision support.}
NASA's concept of operations for an exploration clinical decision support system (CDSS) describes a tool that helps the crew assess and diagnose conditions and decide responses when communication with the ground is delayed or absent, and that eases delayed communication with flight surgeons~\cite{krihak2021cdss}; a later analysis of preliminary data suggests such a system could improve mission outcomes~\cite{russell2023cdss}. Google and NASA have tested a proof-of-concept CMO digital assistant that suggests diagnoses from reported symptoms, to support the CMO or the flight surgeon~\cite{googleCmoDa2025}. The CSA's Connected Care Medical Module (C2M2) programme funds prototype medical modules built around an AI-based system that connects medical devices~\cite{csaC2M2}. A NASA workshop on artificial intelligence for space health recommended an automated, autonomous ``precision space health'' system to monitor and assess crew health~\cite{scott2023}. LIGHT-TIME is narrower than any of these. It makes no diagnosis and uses no trained model to decide anything; it watches a few vital signs, states a deadline and names who can meet it.

\paragraph{Wearables in flight and on Earth.}
The CSA Bio-Monitor, a smart shirt with a headband and a tablet application, records heart activity, blood pressure, breathing, skin temperature, blood oxygen saturation and activity on the ISS, and sends the data to the ground~\cite{csaBioMonitor}. LIGHT-TIME assumes a garment of this kind as its sensor. On Earth, smartwatch heart rate and step counts have flagged infection before symptoms by comparing each person with their own baseline~\cite{mishra2020}, and a real-time alerting system built on that idea ran in a prospective study~\cite{alavi2022}. The personal usual range in LIGHT-TIME borrows the idea of a per-person baseline, but shows it as context and raises no alarm from it (Section~\ref{sec:methods}).

\paragraph{Alarm design.}
In intensive care, most monitor alarms are false or not acted on~\cite{drew2014alarm,gorges2009alarm}, and in one unit a short delay before each alarm was presented would have removed many of the ones that staff ignored~\cite{gorges2009alarm}. LIGHT-TIME applies the same idea to a crew's monitor: a band counts only after it has held for \numSustainN{} readings, and nothing can alarm while the sensor settles after a dropout.

\paragraph{Operations under delay.}
NASA's review of analog missions under communication delay recommended visual indicators of the delay, such as a countdown until a message is received, and procedures that rely on onboard tools rather than contact with the ground~\cite{rader2013delay}. LIGHT-TIME shows every message with its time to landing and its earliest reply. A simulator study of crews managing medical events found that crews without a flight surgeon performed worse, and concluded that checklists alone are not enough for autonomous care~\cite{yule2023crew}. The delay is also not the only gap. Near solar conjunction, communication with a Mars crew can be disrupted for days to months~\cite{mcbrayer2022delay}, so LIGHT-TIME has a no-link mode in which the crew act alone and messages wait in a queue.

\begin{table*}[t]\centering\footnotesize
\caption{Systems nearest to LIGHT-TIME, as described in the cited source. ``n.d.'': not described in that source. Only systems whose sources we could open and check are listed.}\label{tab:related}
\begin{tabular}{@{}L{0.15\textwidth}L{0.13\textwidth}L{0.18\textwidth}L{0.11\textwidth}L{0.19\textwidth}L{0.11\textwidth}@{}}\toprule
System & What it is & Earth delay in the logic & Time left to a threshold & Role & Code public \\ \midrule
ExMC CDSS~\cite{krihak2021cdss,russell2023cdss} & NASA concept of operations & delay and loss of contact drive the design & n.d. & helps the crew assess, diagnose and decide & n.d. \\
CMO digital assistant~\cite{googleCmoDa2025} & proof of concept & named as the motivation & n.d. & supports the CMO or flight surgeon from reported symptoms & n.d. \\
C2M2~\cite{csaC2M2} & funded prototypes & autonomy far from Earth is the aim & n.d. & AI-based system connecting medical devices & n.d. \\
Bio-Monitor~\cite{csaBioMonitor} & smart shirt flown on the ISS & n.d. & n.d. & records vital signs and sends them to the ground & n.d. \\
Autonomy levels~\cite{cypko2025} & position paper & four levels by distance and latency & n.d. & framework, not software & n.d. \\
Smartwatch alerting~\cite{alavi2022} & prospective study on Earth & not applicable & n.d. & alerts on deviation from a personal baseline & open-source platform stated \\
LIGHT-TIME (this paper) & hackathon build & light time for the mission date sets who acts & yes, compared with the round trip & rules and statistics; advises, never decides & repository; public release pending \\
\bottomrule
\end{tabular}
\end{table*}
